\section{Past International Actions and Current Mechanisms}

\subsection{UNESCO Initiatives}

UNESCO's work in protecting cultural heritage consists of an international legal structure that puts in place a basic set of standards and provides the means for states to work together. The most prominent is the 1954 Hague Convention for the Protection of Cultural Property in the Event of Armed Conflict, which established the first comprehensive framework for safeguarding cultural property during times of war. Adopted in response to the wide-scale devastation caused by World War II, its primary purpose was to provide legal protection for cultural property such as monuments, museums, and archives. The Convention establishes a system where cultural sites are identified through the use of the Blue Shield emblem and holds states accountable by requiring them to prevent damage to or theft and illegal trade of cultural objects.

UNESCO further strengthened the international heritage protection framework through the 1972 Convention concerning the Protection of the World Cultural and Natural Heritage, as it legally binds member nations to identify, protect, preserve, and transmit cultural and natural heritage sites of "Outstanding Universal Value" for future generations. The Convention for the Safeguarding of the Intangible Cultural Heritage is another binding international treaty designed to protect living cultural traditions with the goals of safeguarding these practices, ensuring respect for the communities that create them, and raising global awareness.

In order to fulfil and implement the obligations stated in these legal agreements, UNESCO has established different operational mechanisms that monitor the condition of cultural heritage sites and consequently help in coordinating the response of the international community should these sites be under threat. These include the World Heritage List, which recognizes cultural and natural properties of Outstanding Universal Value, and the List of World Heritage in Danger, which identifies sites under threat from armed conflict, natural disasters, climate change, and other factors.

UNESCO works with its Member States to assess the state of conservation of World Heritage properties and make recommendations for proper protection, through its various Reactive Monitoring missions and Periodic Reporting. In cases of armed conflict, UNESCO complements its legal framework through operational measures that are intended to prevent the destruction of cultural heritage and support post-conflict recovery. It coordinates emergency response missions, conducts assessments of damage, coordinates with national authorities and heritage professionals to document losses, secure endangered collections, and develop strategies for reconstruction.

Moreover, UNESCO promotes international awareness through campaigns such as \#Unite4Heritage. The campaign was launched in response to the destruction of several cultural heritage sites in Syria, Iraq, and other countries by extremist groups. It encourages governments, cultural institutions, and the public to work together to prevent cultural heritage destruction and illegal trafficking.

UNESCO cooperates closely with Member States to implement practical measures for the protection of cultural heritage, working directly with governments, researchers, cultural institutions, municipalities, and other stakeholders to provide training programmes and resources to site managers, communities, and policymakers. It focuses its contributions on practical workshops, disaster management, and digital documentation to sustainably safeguard tangible heritage.

Key UNESCO initiatives include the World Heritage Capacity Building Strategy, which promotes training, institutional development, and knowledge-sharing among heritage professionals and local communities, Priority Africa (Flagship Programme 3), targeting World Heritage conservation particularly in African states, strengthening their capacities to fight against illicit trafficking of cultural property and the preservation and promotion of collections and museums, and many other initiatives supporting specifically the protection of intangible heritage.

In each of these areas, UNESCO's actions are carried out in collaboration with states with the consent and participation of the national authorities concerned. UNESCO provides its expertise, methodological tools, and platforms on which to exchange ideas, while the national and local authorities have the main role as far as the development and implementation of any necessary policies and measures are concerned. Beyond direct cooperation with Member States, UNESCO also promotes international cooperation. Through global and regional platforms, the organization facilitates discussions about policy development, best practices, and funding priorities and leads global debates on heritage and intercultural dialogue, providing space for nations and experts to share their experiences and struggles.

\subsection{International Law and the International Criminal Court}

\subsubsection[International Humanitarian Law and Cultural Property]{International Humanitarian Law and Cultural Property\footnote{International Committee of the Red Cross. (2005). \emph{Customary international humanitarian law, Volume I: Rules}. Cambridge University Press.\url{https://www.icrc.org/en/doc/assets/files/other/customary-international-humanitarian-law-i-icrc-eng.pdf}f}}

Three foundational IHL principles are directly relevant:

\begin{itemize}
  \item distinction = requires parties to distinguish between civilian objects and military objectives;
  \item proportionality = prohibits attacks causing excessive incidental damage; and
  \item military necessity = permits only measures which are indispensable to achieving a legitimate military objective.
\end{itemize}

Cultural property qualifies as a civilian object under IHL unless used for military purposes, a condition frequently invoked by parties seeking to justify attacks on heritage sites.\footnote{Henckaerts, J. M., \& Doswald-Beck, L. (2005). \emph{Customary international humanitarian law, Volume I: Rules}. Cambridge University Press / ICRC.\url{https://doi.org/10.1017/CBO9780511804700}}

Building on these principles, the First Protocol of the 1954 Hague Convention obliged occupying powers to prevent the export of cultural property from occupied territories, while the Second Protocol of 1999 introduced enhanced protection for the most significant sites and, critically, established individual criminal responsibility for serious violations, the first time the Hague framework incorporated individual accountability rather than relying exclusively on state responsibility.\footnote{UNESCO. (1954). \emph{Convention for the Protection of Cultural Property in the Event of Armed Conflict}.\url{https://en.unesco.org/protecting-heritage/convention-and-protocols/1954-convention}} \footnote{UNESCO. (1999). \emph{Second Protocol to the Hague Convention of 1954 for the Protection of Cultural Property in the Event of Armed Conflict}.\url{https://en.unesco.org/protecting-heritage/convention-and-protocols/second-protocol}} As of 2024, however, only 85 states had ratified the Second Protocol, leaving a significant portion of the international community outside its enhanced protection regime.\footnote{UNESCO. (2024). \emph{States parties to the Second Protocol}.\url{https://en.unesco.org/protecting-heritage/convention-and-protocols/states-parties}}

\subsubsection{The Rome Statute and Individual Criminal Responsibility}

The Rome Statute of the ICC, entering into force in 2002, provides the principal mechanism for individual criminal accountability. Article 8(2)(b)(ix) criminalises intentional attacks against historic monuments and buildings dedicated to religion, education, art or science in international armed conflict, while Article 8(2)(e)(iv) replicates this prohibition for non-international armed conflict, a significant inclusion given that the majority of heritage destruction since 1998 has occurred in internal conflicts involving non-state actors.\footnote{International Criminal Court. (1998). \emph{Rome Statute of the International Criminal Court}.\url{https://www.icc-cpi.int/resource-library/documents/rs-eng.pdf}}

Nevertheless, the Court may only exercise jurisdiction over nationals of states parties or over crimes committed on state party territory, unless the Security Council refers a situation under Chapter VII of the UN Charter meaning nationals of non-signatory states are effectively beyond the Court's reach without a referral that the veto structure makes nearly impossible when a permanent member is involved.\footnote{Francioni, F. (2004). Beyond state sovereignty: The protection of cultural heritage as a shared interest of humanity. \emph{Michigan Journal of International Law, 25}(4), 1209–1228.}

\subsubsection{The Al Mahdi Case: Legal Precedent and Remaining Questions}

The 2016 conviction of Ahmad Al Faqi Al Mahdi, produced three legally significant outcomes. The Trial Chamber confirmed that intentional destruction of cultural sites constitutes a war crime even without direct loss of human life; the reparations order established that victims are entitled to collective reparations including reconstruction costs; and the case demonstrated that prosecution of non-state actors for heritage destruction is legally feasible, even if the conditions that made it possible remain exceptional.\footnote{International Criminal Court. (2016). \emph{Prosecutor v. Ahmad Al Faqi Al Mahdi, Judgment and Sentence} (ICC-01/12-01/15).\url{https://www.icc-cpi.int/mali/al-mahdi}} \footnote{International Criminal Court. (2017). \emph{Prosecutor v. Ahmad Al Faqi Al Mahdi, Reparations Order} (ICC-01/12-01/15-236).\url{https://www.icc-cpi.int/mali/al-mahdi}}

The case did not address how the Rome Statute applies to heritage destruction carried out by state armed forces, nor did it resolve the question of command responsibility, the doctrine under which military commanders may be held responsible for crimes committed by forces under their effective control. These are gaps directly relevant to ongoing situations where state or state-affiliated actors are responsible for documented heritage destruction.\footnote{Schabas, W. A. (2016). \emph{The International Criminal Court: A commentary on the Rome Statute} (2nd ed.). Oxford University Press.}

\subsubsection{State Responsibility}

Where the Rome Statute cannot reach, the International Law Commission's Articles on Responsibility of States for Internationally Wrongful Acts (2001) establish that a state bears international responsibility for breaches of obligations under the 1954 Hague Convention and customary IHL attributable to it.\footnote{International Law Commission. (2001). \emph{Articles on Responsibility of States for Internationally Wrongful Acts}. United Nations.\url{https://legal.un.org/ilc/texts/instruments/english/draft_articles/9_6_2001.pdf}}

In principle, such a state could be brought before the International Court of Justice. In practice, no state has ever been successfully held responsible before an international court specifically for heritage destruction, a gap reflecting both political constraints on interstate litigation and the structural limitation of a system entirely dependent on states choosing to bring other states to account.\footnote{Forrest, C. (2010). \emph{International law and the protection of cultural heritage}. Routledge.}

\subsection{UN Resolutions and International Cooperation}

\textbf{Resolution 2347 (2017)} was the first Security Council resolution dedicated exclusively to cultural heritage protection, formally recognising intentional heritage destruction as a tactic of war and a potential threat to international peace and security. It called upon states to develop national inventories of cultural property, strengthen border controls, improve judicial cooperation and support UNESCO's monitoring capacity.\footnote{United Nations Security Council. (2017). \emph{Resolution 2347} (S/RES/2347).\url{https://undocs.org/S/RES/2347(2017)}}

\textbf{Resolution 2199 (2015)} condemned ISIS's destruction and prohibited trade in illegally removed cultural property from Iraq and Syria.\footnote{United Nations Security Council. (2015). \emph{Resolution 2199} (S/RES/2199).\url{https://undocs.org/S/RES/2199(2015)}}

\textbf{Resolution 2509 (2020)} and \textbf{Resolution 2585 (2021)} addressed heritage protection in Libya and Syria respectively, urging compliance with IHL obligations and calling on Member States to prevent illicit trafficking from those territories.\footnote{United Nations Security Council. (2020). \emph{Resolution 2509} (S/RES/2509).\url{https://undocs.org/S/RES/2509(2020)}} \footnote{United Nations Security Council. (2021). \emph{Resolution 2585} (S/RES/2585).\url{https://undocs.org/S/RES/2585(2021)}}

\textbf{General Assembly Resolution 69/281 (2015)} called for a global response to deliberate heritage destruction and illicit trafficking of cultural property and consistently referenced heritage protection as a component of post-conflict reconciliation.\footnote{United Nations General Assembly. (2015). \emph{Resolution 69/281} (A/RES/69/281).\url{https://undocs.org/A/RES/69/281}}

International cooperation has proven equally essential in translating these normative commitments into practice; INTERPOL's Works of Art unit coordinates cross-border law enforcement efforts to identify and recover looted cultural property;\footnote{INTERPOL. (2022). \emph{Works of art: Protecting cultural heritage}.\url{https://www.interpol.int/en/Crimes/Cultural-heritage-crime}} the World Customs Organization has developed specific guidelines for customs officials to intercept illicitly trafficked cultural goods at border crossings;\footnote{World Customs Organization. (2020). \emph{Guidelines on combating customs offences relating to cultural property}.\url{https://www.wcoomd.org/}} and the European Union adopted Council Regulation 2019/880, establishing harmonised import controls specifically targeting cultural property removed from conflict zones.\footnote{European Union. (2019). \emph{Council Regulation (EU) 2019/880 on the introduction and the import of cultural goods}.\url{https://eur-lex.europa.eu/legal-content/EN/TXT/?uri=CELEX:32019R0880}}

However, the current system faces structural limitations: Security Council resolutions on cultural heritage are adopted under Chapter VI rather than Chapter VII of the UN Charter, meaning they call upon rather than compel states to act, with no automatic enforcement consequences for non-compliance;\footnote{Forrest, C. (2010). \emph{International law and the protection of cultural heritage}. Routledge.} the definition of cultural heritage itself remains contested, with significant disagreements over what qualifies as heritage worthy of protection and whose heritage counts; and state sovereignty continues to be routinely invoked to resist external scrutiny of domestic heritage policy.\footnote{Blake, J. (2000). On defining the cultural heritage. \emph{International and Comparative Law Quarterly, 49}(1), 61–85.\url{https://doi.org/10.1017/S002058930006396X}}
